The goal of this chapter is to present some polynomial-time algorithms
for convex optimization. These algorithms use knowledge of the function
in a minimal way, essentially by querying the function value along
with weak bounds on its support and range. Historically, the ideas
in this chapter originated in the work of \cite{Grunbaum1960,yudin1976evaluation},
developed into polynomial-time algorithms by \cite{khachiyan1980polynomial,grotschel1981ellipsoid,bertsimas2004solving}.
For a comprehensive presentation of related problems and ellipsoid-based
algorithms, the reader is referred to the classic book \cite{grotschel2012geometric}. 

\section{Cutting Plane Methods\label{sec:Cutting-Plane-Methods}}

Given a continuously differentiable convex function $f$, Theorem
\ref{thm:first_order} shows that
\begin{equation}
f(y)\geq f(x)+\left\langle \nabla f(x),y-x\right\rangle \text{ for all }y.\label{eq:sub_grad_def}
\end{equation}
Let $x^{*}$ be any minimizer of $f$. Replacing $y$ with $x^{*}$,
we have that 
\[
f(x)\geq f(x^{*})\geq f(x)+\left\langle \nabla f(x),x^{*}-x\right\rangle .
\]
Therefore, we know that $\left\langle \nabla f(x),x^{*}-x\right\rangle \leq0$.
Namely, $x^{*}$ lies in a halfspace $H$ with normal vector $-\nabla f(x)$.
Roughly speaking, this shows that each gradient computation cuts the
set of possible solutions in ``half''. In one dimension, this allows
us to do a binary search to minimize convex functions. 

It turns out that in $\Rn$, binary search still works. In this chapter,
we will cover several ways to do this binary search. All of them follow
the same framework, called the \emph{cutting plane method}. In this
method, the convex set or function of interest is given by an oracle,
typically a \emph{separation oracle} (for definitions in a broader
context, see Defs. \ref{def:SEP} and \ref{def:GRAD} in the next
chapter): for any $x\notin K\subseteq\R^{n}$, the oracle finds a
vector $g(x)\in\R^{n}$ such that
\[
g(x)^{\top}(y-x)\leq0\text{ for all }y\in K.
\]

Cutting plane methods address the following class of problems.
\begin{problem}[Finding a point in a convex set]
\label{prob:abstract_convex_problem}Given $\epsilon>0,R>0,$ and
a convex set $K\subseteq R\cdot B^{n}$ specified by a separation
oracle, find a point $y\in K$ or conclude that $\vol(K)\leq\epsilon^{n}$.
The complexity of an algorithm is measured by the number of calls
to the oracle and the number of arithmetic operations.
\end{problem}

\begin{rem*}
To minimize a convex function, we set $g(x)=\nabla f(x)$ and $K$
to be the set of (approximate) minimizers of $f$. In Section \ref{sec:vol2val},
we relate the problem of proving that $\vol(K)$ is small to the problem
of finding an approximate minimizer of $f$.
\end{rem*}
In this framework, we maintain a convex set $E^{(k)}$ that contains
the set $K$. In each iteration, we choose some $x^{(k)}$ based on
$E^{(k)}$ and query the oracle for $g(x^{(k)})$ . The guarantee
for $g(.)$ implies that $K$ lies in the halfspace 
\[
H^{(k)}=\{y:g(x^{(k)})^{\top}(y-x^{(k)})\leq0\}
\]
 and hence $K\subset H^{(k)}\cap E^{(k)}.$ The algorithm continues
by choosing $E^{(k+1)}$ to be a convex set that contains $H^{(k)}\cap E^{(k)}$. 

\begin{algorithm2e}[H]

\caption{$\mathtt{CuttingPlaneFramework}$}\label{alg:cutting_plane}

\SetAlgoLined

\textbf{Input:} Initial set $E^{(0)}\subseteq\Rn$ containing $K$.

\For{$k=0,1,\cdots$}{

Choose a point $x^{(k)}\in E^{(k)}$.

\lIf{$E^{(k)}$ is ``small enough''}{\Return$x^{(k)}$}

Find $E^{(k+1)}\supset E^{(k)}\cap H^{(k)}$ where
\begin{equation}
H^{(k)}\defeq\{x\in\Rn\text{ such that }g(x^{(k)})^{\top}(x-x^{(k)})\leq0\}.\label{eq:Hk}
\end{equation}

}

\end{algorithm2e}

To analyze the algorithm, the main questions we need to answer are:
\begin{enumerate}
\item How do we choose $x^{(k)}$ and $E^{(k+1)}$?
\item How do we measure progress?
\item How quickly does the method converge?
\item How expensive is each step?
\end{enumerate}
Progress on the cutting plane method is shown in the next table.

\begin{table}[h]
\begin{centering}
{\footnotesize{}%
\begin{tabular}{|c|c|c|c|}
\hline 
{\footnotesize Year} & {\footnotesize$E^{(k)}$ and $x^{(k)}$} & {\footnotesize Iter} & {\footnotesize Cost/Iter}\tabularnewline
\hline 
\hline 
{\footnotesize 1965 \cite{levin1965algorithm,newman1965location}} & {\footnotesize Center of gravity} & {\footnotesize$n$} & {\footnotesize$n^{n}$}\tabularnewline
\hline 
{\footnotesize 1979 \cite{yudin1976evaluation,shor1977cut,khachiyan1980polynomial}} & {\footnotesize Center of ellipsoid} & {\footnotesize$n^{2}$} & {\footnotesize$n^{2}$}\tabularnewline
\hline 
{\footnotesize 1988 \cite{khachiyan1988inscribed}} & {\footnotesize Center of John ellipsoid} & {\footnotesize$n$} & {\footnotesize$n^{2.878}$}\tabularnewline
\hline 
{\footnotesize 1989 \cite{DBLP:conf/focs/Vaidya89b}} & {\footnotesize Volumetric center} & {\footnotesize$n$} & {\footnotesize$n^{2.378}$}\tabularnewline
\hline 
{\footnotesize 1995 \cite{atkinson1995cutting}} & {\footnotesize Analytic center} & {\footnotesize$n$} & {\footnotesize$n^{2.378}$}\tabularnewline
\hline 
{\footnotesize 2004 \cite{bertsimas2004solving}} & {\footnotesize Center of gravity} & {\footnotesize$n$} & {\footnotesize$n^{4}$}\tabularnewline
\hline 
{\footnotesize 2015 \cite{lee2015faster}} & {\footnotesize Hybrid center} & {\footnotesize$n$} & {\footnotesize$n^{2}$ (amortized)}\tabularnewline
\hline 
{\footnotesize 2020 \cite{jiang2020improved}} & {\footnotesize Volumetric center} & {\footnotesize$n$} & {\footnotesize$n^{2}$ (amortized)}\tabularnewline
\hline 
\end{tabular}}{\footnotesize\par}
\par\end{centering}
\centering{}{\footnotesize\caption{{\footnotesize\label{table:cpm}Different Cutting Plane Methods. Omitted
polylogarithmic terms. The number of iterations follows from the rate.}}
}{\footnotesize\par}
\end{table}


\section{Ellipsoid Method}

\begin{figure}[h]
\begin{centering}
\includegraphics[width=0.75\columnwidth]{Figures/Section_3\lyxdot 2}
\par\end{centering}
\caption{Ellipsoid Method}

\end{figure}

We start by explaining the Ellipsoid method. The algorithm maintains
an ellipsoid
\[
E^{(k)}\defeq\{y\in\Rn:\ (y-x^{(k)})^{\top}\left(A^{(k)}\right)^{-1}(y-x^{(k)})\leq1\}
\]
that contains $K$ and becomes smaller in volume in each step. Note
that for this to be an ellipsoid the matrix $A^{(k)}$ must be symmetric
PSD. After we compute $g(x^{(k)})$ and $H^{(k)}$ via (\ref{eq:Hk}),
we define $E^{(k+1)}$ to be the smallest volume ellipsoid containing
$E^{(k)}\cap H^{(k)}$. The key observation is that the volume of
the ellipsoid $E^{(k)}$ decreases by a factor of $1-\Theta(\frac{1}{n})$
every iteration. This volume property holds for any halfspace through
the center of the current ellipsoid (not only for the one whose normal
is the gradient), a property we will exploit in the next chapter.

\begin{algorithm2e}[H]

\caption{$\mathtt{Ellipsoid}$}\label{alg:ellipsoid_method}

\SetAlgoLined

\textbf{Input:} Initial ellipsoid $E^{(0)}=\{y\in\Rn:\ (y-x^{(0)})^{\top}\left(A^{(0)}\right)^{-1}(y-x^{(0)})\leq1\}$.

\For{$k=0,1,\cdots$}{

\lIf{$E^{(k)}$ is ``small enough'' or Oracle says YES}{\Return$x^{(k)}$}

$x^{(k+1)}=x^{(k)}-\frac{1}{n+1}\frac{A^{(k)}g(x^{(k)})}{\sqrt{g(x^{(k)})^{\top}A^{(k)}g(x^{(k)})}}$.

$A^{(k+1)}=\frac{n^{2}}{n^{2}-1}\left(A^{(k)}-\frac{2}{n+1}\frac{A^{(k)}g(x^{(k)})g(x^{(k)})^{\top}A^{(k)}}{g(x^{(k)})^{\top}A^{(k)}g(x^{(k)})}\right)$.

}

\end{algorithm2e}
\begin{rem}
We usually obtain $g(x)$ from a \emph{separation oracle}.
\end{rem}

\begin{lem}
\label{lem:ellipsoid_vol_decrease}For the Ellipsoid method (Algorithm
\ref{alg:ellipsoid_method}), we have 
\[
\vol E^{(k+1)}<e^{-\frac{1}{2n+2}}\vol E^{(k)}\quad\mbox{and}\quad E^{(k)}\cap H^{(k)}\subset E^{(k+1)}.
\]
\end{lem}

\begin{rem}
Note that the proof below also shows that $\vol E^{(k+1)}=e^{-\Theta(\frac{1}{n})}\vol E^{(k)}$.
Therefore, the ellipsoid method does not run faster even for any function,
e.g., a linear function, making it provably ``slow'' in practice.
\end{rem}

\begin{proof}
Note that the ratio of $\vol E^{(k+1)}/\vol E^{(k)}$ does not change
under any affine transformation, and neither does set containment.
Therefore, we can do a transformation so that $A^{(k)}=I$, $x^{(k)}=0$
and $v(x^{(k)})=e_{1}$. Therefore $H^{(k)}$ is simply the halfspace
$\{x:x_{1}\le0\}$. This simplifies our calculation. We need to prove
two statements: $\vol E^{(k+1)}<e^{-\frac{1}{2n+2}}\vol E^{(k)}$
and that $E^{(k)}\cap H^{(k)}\subset E^{(k+1)}$.

Claim 1: $\vol E^{(k+1)}<e^{-\frac{1}{2(n+1)}}\vol E^{(k)}$.

Note that since $v(x^{(k)})=e_{1},$we have $A^{(k+1)}=\frac{n^{2}}{n^{2}-1}\left(I-\frac{2}{n+1}e_{1}e_{1}^{\top}\right)$.
Therefore, we have that
\begin{align*}
\left(\frac{\vol E^{(k+1)}}{\vol E^{(k)}}\right)^{2}=\left|\frac{\det A^{(k+1)}}{\det A^{(k)}}\right| & =\left(\frac{n^{2}}{n^{2}-1}\right)^{n}\det\left(I-\frac{2}{n+1}e_{1}e_{1}^{\top}\right)\\
 & =\left(\frac{n^{2}}{n^{2}-1}\right)^{n-1}\frac{n^{2}}{n^{2}-1}\left(\frac{n-1}{n+1}\right)\\
 & =\left(1+\frac{1}{n^{2}-1}\right)^{n-1}\left(1-\frac{1}{n+1}\right)^{2}\\
 & <\exp(\frac{n-1}{n^{2}-1}-\frac{2}{n+1})=\exp(-\frac{1}{n+1})
\end{align*}
where we used $1+x\leq e^{x}$ for all $x$.

Claim 2. $E^{(k)}\cap H^{(k)}\subset E^{(k+1)}$.

By the definition of $E^{(k)}$and the assumption $A^{(k)}=I$, we
have 
\[
x^{(k+1)}=-\frac{e_{1}}{n+1}
\]
and, for any $x\in E^{(k)}\cap H^{(k)}$, we have that $\|x\|_{2}\leq1$
and $x_{1}\leq0$. By direct computation, using the fact that 
\[
\left(I-\frac{2}{n+1}e_{1}e_{1}^{\top}\right)^{-1}=\left(I+\frac{2}{n-1}e_{1}e_{1}^{\top}\right)
\]
we have
\begin{align*}
 & (x+\frac{e_{1}}{n+1})^{\top}\left(1-\frac{1}{n^{2}}\right)\left(I+\frac{2}{n-1}e_{1}e_{1}^{\top}\right)(x+\frac{e_{1}}{n+1})\\
= & \left(1-\frac{1}{n^{2}}\right)\left(\|x\|^{2}+\frac{2x_{1}(1+x_{1})}{n-1}+\frac{1}{n^{2}-1}\right)\\
\leq & \left(1-\frac{1}{n^{2}}\right)\left(1+0+\frac{1}{n^{2}-1}\right)=1
\end{align*}
where we used that $\|x\|_{2}\leq1$ and $x_{1}(1+x_{1})\leq0$ (since
$-1\leq x_{1}\leq0$) at the end. This shows that $x\in E^{(k+1)}$.
\end{proof}
\begin{xca}
Show that the ellipsoid $E^{(k+1)}$ computed above is the minimum
volume ellipsoid containing $E^{(k)}\cap H^{(k)}$.
\end{xca}

\begin{xca}
Suppose that we used a box instead of an ellipsoid. Could we ensure
progress in each iteration? What about a simplex?
\end{xca}


\section{From Volume to Function Value\label{sec:vol2val}}

Lemma \ref{lem:ellipsoid_vol_decrease} shows that the volume of a
set containing all minimizers decreases by a constant factor every
$n$ steps. In general, knowing that the optimal $x^{*}$ lies in
a small volume set does not provide enough information to find a point
with small function value. For example, if we only knew that $x^{*}$
lies in the plane $\{x:x_{1}=0\}$, we would still need to search
for $x^{*}$ over an $n-1$ dimensional space. However, if the set
is constructed by the cutting plane framework (Algorithm \ref{alg:cutting_plane}),
then we can guarantee that small volume implies that any point in
the set has close to optimal function value.

This implies that we can minimize any convex function with $\varepsilon$
additive error in $O(n^{2}\log(1/\varepsilon))$ iterations. To make
the statement more general, we note that the Ellipsoid method can
be used for non-differentiable functions. In the next theorem, we
allow for a general ``progress'' function and an arbitrary support
$\Omega$. The convergence argument works for any monotonic, scale-and-shift-invariant
progress measure. To allow for non-differentiable convex functions,
in place of gradient, we use the subgradient (see Sec. \ref{sec:Subgradients}). 
\begin{thm}
\label{thm:conv_opt}Let $x^{(k)}$ be the sequence of points produced
by the cutting plane framework (Algorithm \ref{alg:cutting_plane})
for a convex function $f$. Let $\mathcal{V}$ be a mapping from subsets
of $\Rn$ to non-negative real numbers satisfying
\begin{enumerate}
\item (Linearity) For any set $S\subseteq\R^{n}$, any vector $y$ and any
scalar $\alpha\geq0$, we have $\mathcal{V}(\alpha S+y)=\alpha\mathcal{V}(S)$
where $\alpha S+y=\{\alpha x+y:x\in S\}$.
\item (Monotonicity) For any set $T\subset S$, we have that $\mathcal{V}(T)\leq\mathcal{V}(S)$.
\end{enumerate}
Then, for any set $\Omega\subseteq E^{(0)}$, with $\mathcal{V}(\Omega)>0$,
and $k$ large enough so that $\mathcal{V}(E^{(k)})\le\mathcal{V}(\Omega)$,
we have
\[
\min_{i=1,2,\cdots k}f(x^{(i)})-\min_{y\in\Omega}f(y)\leq\frac{\mathcal{V}(E^{(k)})}{\mathcal{V}(\Omega)}\cdot\left(\max_{z\in\Omega}f(z)-\min_{x\in\Omega}f(x)\right)\enspace.
\]
\end{thm}

\begin{rem}
We can think of $\mathcal{V}(E)$ as some way to measure the size
of $E$. It can be radius, mean-width or any other way to measure
``size''. For the ellipsoid method, we use $\mathcal{V}(E)=\vol(E)^{\frac{1}{n}}$
for which we have proved volume decrease in Lemma \ref{lem:ellipsoid_vol_decrease}.
We raise the volume to power $1/n$ to satisfy linearity. Also note
that we only guarantee that one of the previous query points has a
small function value; of course we can simply use the point with minimum
function value.
\end{rem}

\begin{proof}
Let $x^{*}$ be any minimizer of $f$ over $\Omega$. For any $1\ge\alpha>\frac{\mathcal{V}(E^{(k)})}{\mathcal{V}(\Omega)}$
and $S=(1-\alpha)x^{*}+\alpha\Omega$, by the linearity of $\mathcal{V}$,
we have that 
\[
\mathcal{V}(S)=\alpha\mathcal{V}(\Omega)>\mathcal{V}(E^{(k)})
\]
Therefore, $S$ is not a subset of $E^{(k)}$ and hence there is a
point $y\in S\backslash E^{(k)}$. $y$ is not in $E^{(k)}$. This
means it is eliminated by the subgradient halfspace at some step $i\leq k$,
namely for some $i\le k$, we have (denoting the subgradient by $\nabla f$),
\[
\nabla f(x^{(i)})^{\top}(y-x^{(i)})>0.
\]
By the convexity of $f$, it follows that $f(x^{(i)})\leq f(y).$
Since $y\in S$, we have $y=(1-\alpha)x^{*}+\alpha z$ for some $z\in\Omega$.
Thus, the convexity of $f$ implies that
\[
f(x^{(i)})\leq f(y)\leq(1-\alpha)f(x^{*})+\alpha f(z).
\]
Therefore, we have
\[
\min_{i=1,2,\cdots k}f(x^{(i)})-\min_{x\in\Omega}f(x)\leq\alpha\left(\max_{z\in\Omega}f(z)-\min_{x\in\Omega}f(x)\right).
\]
Since this holds for any $\alpha>\frac{\mathcal{V}(E^{(k)})}{\mathcal{V}(\Omega)}$,
we have the result.
\end{proof}
Combining Lemma \ref{lem:ellipsoid_vol_decrease} and Theorem \ref{thm:conv_opt},
we have the following rate of convergence.
\begin{thm}
\label{thm:ellipsoid}Let $f$ be a convex function on $\Rn$, $E^{(0)}$
be any initial ellipsoid and $\Omega\subset E^{(0)}$ be any convex
set. Suppose that for any $x\in E^{(0)}$, we can find, in time $\mathcal{T}$,
a nonzero vector $g(x)$ such that 
\[
f(y)\geq f(x)\text{ for any }y\text{ such that }g(x)^{\top}(y-x)\geq0.
\]

Then, we have
\[
\min_{i=1,2,\cdots k}f(x^{(i)})-\min_{y\in\Omega}f(y)\leq\left(\frac{\vol(E^{(0)})}{\vol(\Omega)}\right)^{\frac{1}{n}}\exp\left(-\frac{k}{2n(n+1)}\right)\left(\max_{z\in\Omega}f(z)-\min_{x\in\Omega}f(x)\right)\enspace.
\]
Moreover, each iteration takes $O(n^{2}+\mathcal{T})$ time.
\end{thm}

\begin{rem}
We note that for this rate of convergence (and hence the entire algorithm)
to be polynomial, we need some bound on the range of the function
value. This often follows from a bound on the diameter of the support.
\end{rem}

\begin{proof}
Lemma \ref{lem:ellipsoid_vol_decrease} shows that the volume of the
ellipsoid maintained decreases by a factor of $\exp(-\frac{1}{2n+2})$
in every iteration. Hence, $\vol^{\frac{1}{n}}$ decreases by $\exp(-\frac{1}{2n(n+1)})$
every iteration. The bound follows by applying Theorem \ref{thm:conv_opt}
with $\mathcal{V}(E)\defeq\vol(E)^{\frac{1}{n}}$. Using Theorem \ref{thm:conv_opt},
we have 
\begin{align*}
\min_{i=1,2,\cdots k}f(x^{(i)})-\min_{y\in\Omega}f(y) & \le\left(\frac{\vol(E^{(k)})}{\vol(\Omega)}\right)^{\frac{1}{n}}\left(\max_{z\in\Omega}f(z)-\min_{x\in\Omega}f(x)\right)\\
 & \le\left(\frac{\vol(E^{(0)})}{\vol(\Omega)}\right)^{\frac{1}{n}}\exp\left(-\frac{k}{2n(n+1)}\right)\left(\max_{z\in\Omega}f(z)-\min_{x\in\Omega}f(x)\right).
\end{align*}

Next, we note the proof of Theorem \ref{thm:conv_opt} only used the
fact that one side of hyperplane defined by the gradient has higher
value. Therefore, we can replace the gradient with the vector $g(x)$. 

To bound the time per iteration, note that we make one query to the
separation oracle per iteration, and then compute the next Ellipsoid
using the formulas in Algorithm \ref{alg:ellipsoid_method}. The most
time-consuming operation is multiplying an $n\times n$ matrix by
an $n$-vector, which has complexity $O(n^{2})$.
\end{proof}
This theorem can be used to solve many problems in polynomial time.
As an illustration, we show how to solve linear programs to desired
approximation in polynomial time here. 
\begin{thm}
Given a linear program $\min_{x\in P}c^{\top}x$ where $P=\{x\text{: }Ax\geq b\}$.
Let the diameter of $P$ be $R\defeq\max_{x\in P}\norm x_{2}$ and
its volume radius be $r=\vol(P)^{1/n}$. Then, we can find $x\in P$
for which 
\[
c^{\top}x-\min_{y\in P}c^{\top}y\leq\varepsilon\cdot(\max_{y\in P}c^{\top}y-\min_{y\in P}c^{\top}y)
\]
in $O(n^{2}(n^{2}+\nnz(A))\log(\frac{R}{r\varepsilon}))$ time where
$\nnz(A)$ is the number of non-zero elements in $A$.
\end{thm}

\begin{rem}
If the dimension $n$ is constant, this algorithm is nearly linear
time (linear in the number of constraints)!
\end{rem}

\begin{proof}
For the linear program $\min_{Ax\geq b}c^{\top}x$, the function we
want to minimize is
\begin{equation}
L(x)=c^{\top}x+\delta_{Ax\geq b}(x)\quad\text{where}\quad\delta_{Ax\geq b}(x)=\begin{cases}
0 & \text{if }a_{i}^{\top}x\geq b_{i}\text{ for all }i\\
+\infty & \text{otherwise.}
\end{cases}.\label{eq:LP_func}
\end{equation}

For this function $L$, we can use the separation oracle $v(x)=c$
if $Ax\geq b$ and $v(x)=-a_{i}$ if $a_{i}^{\top}x<b_{i}$. If there
are many violated constraints, any one of them will do.

In this case, we can set $\Omega=P$. We can simply pick $E^{(0)}$
be the unit ball centered at $0$ with radius $R$. We apply Theorem
\ref{thm:ellipsoid} to find $x$ such that
\[
c^{\top}x-\min_{y\in P}c^{\top}y\leq\varepsilon\cdot(\max_{y\in P}c^{\top}y-\min_{y\in P}c^{\top}y)
\]
in time $O(n^{2}(n^{2}+\nnz(A))\log(\frac{R}{r\varepsilon})).$ Note
that the complexity of computing a matrix-vector product $Ax$ is
$O(\nnz(A))$. 
\end{proof}
\begin{xca}
Give a bound on $R$ as a polynomial in $n,\langle A\rangle,\langle b\rangle$,
where $\langle A\rangle,\langle b\rangle$ are the numbers of bits
needed to represent $A,b$ whose entries are rationals.
\end{xca}

To get the solution exactly, i.e., $\varepsilon=0$, we need to assume
the linear program has integral (or rational) coefficients and then
the running time will depend on the sizes of the numbers in the matrix
$A$ and in the vectors $b$ and $c$. It is still open how to solve
linear programs in time bounded by a polynomial in only the number
of variables and constraints (and not the bit sizes of the coefficients).
Such a running time is called \emph{strongly polynomial. }
\begin{problem*}
Can we solve linear programs in strongly polynomial time?
\end{problem*}

\section{Center of Gravity Method}

In the Ellipsoid method, we used an ellipsoid as the current set,
and its center as the next query point.

In the center of gravity method, we start with any bounded convex
set containing a minimizer, e.g., a large enough cube, and the set
maintained is simply the intersection of all halfspaces used so far
with the original set. The query point will be the center of gravity
(or barycenter) of the current set. Recall that the center of gravity
of a bounded set is the average of points in the set, i.e., 
\[
z_{K}=\frac{1}{\vol(K)}\int_{K}x\:dx.
\]
The measure of progress will once again be the volume (or more precisely,
volume radius) of the current set.

It is clear that the volume can only decrease in each iteration. But
at what rate? The following classical theorem shows that the volume
of the convex body decreases by a constant factor (no more than $1-\frac{1}{e}$)
when using the exact center of gravity.
\begin{thm}
\label{thm:Grunbaum}\cite{Grunbaum1960}Let $K$ be a convex body
in $\R^{n}$ with center of gravity $z$. Let $H$ be any halfspace
containing $z$. Then, 
\[
\vol(K\cap H)\ge\left(\frac{n}{n+1}\right)^{n}\vol(K).
\]
\end{thm}

Note that the constant on the RHS is at least $1/e$. We prove the
theorem later in this chapter. Unfortunately, computing the center
of gravity even of a polytope is \#P-hard \cite{rademacher2007approximating}.
For the purpose of efficient approximations, it is important to establish
a stable version of the theorem that does not require an exact center
of gravity.

Recall that a nonnegative function is logconcave if its logarithm
is concave, i.e., for any $x,y\in\R^{n}$ and any $\lambda\in[0,1]$,
\[
f(\lambda x+(1-\lambda)y)\ge f(x)^{\lambda}f(y)^{1-\lambda}.
\]
We refer to Section \ref{sec:Logconcave-functions} for some background
on logconcave functions. A distribution $p$ is \emph{isotropic} if
the mean of a random variable drawn from the distribution is zero
and the covariance is the identity matrix.

The randomized center of gravity method defined as follows: 

\begin{algorithm2e}[H]

\caption{$\mathtt{RandomizedCenterOfGravity}$}\label{alg:randomized_cog}

\SetAlgoLined

\textbf{Input:} Initial convex set $E^{(0)}$.

\For{$k=0,1,\cdots$}{

\lIf{$E^{(k)}$ is ``small enough''}{\Return$x^{(k)}$}

Let $y^{(1)},\ldots,y^{(N)}$ be uniform random points from $E^{(k)}$
and set
\begin{align*}
x^{(k)} & =\frac{1}{N}\sum_{i=1}^{N}y^{(i)}\\
E^{(k+1)} & =E^{(k)}\cap H^{(k)}
\end{align*}

where $H^{(k)}\defeq\{x\in\Rn:g(x^{(k)})^{\top}(x-x^{(k)})\leq0\}$
obtained by querying $g(x^{(k)}).$ }

\end{algorithm2e}
\begin{rem}
The question of how to sample (nearly) uniform random points is an
important one that we will address in detail in this book. Using sampling
to approximate the center of gravity leads to the reduction in the
time per iteration from $n^{n}$ to $n^{4}$ (Table \ref{table:cpm}).
\end{rem}

To prove the convergence of the method, we will use a robust version
of Theorem \ref{thm:Grunbaum}, which will give a similar result despite
using only an approximate center of gravity.
\begin{thm}[Robust Gr\"unbaum]
\label{thm:grun-1}Let $p$ be an isotropic logconcave distribution,
namely $\E_{x\sim p}x=0$ and $\E_{x\sim p}x^{2}=1$. For any unit
vector $\theta\in\R^{n}$, $t\in\R$ we have
\[
\P_{x\sim p}(x^{\top}\theta\geq t)\geq\frac{1}{e}-\left|t\right|.
\]
\end{thm}

\begin{proof}
By taking the marginal with respect to the direction $\theta$, we
can assume the distribution is one-dimensional. Let $P(t)=\P_{x\sim p}(x^{\top}\theta\le t)$.
Note that $P(t)$ is the convolution of $p$ and $1_{(-\infty,0]}$.
Hence, it is logconcave (Lemma \ref{lem:marginal}). By some limit
arguments, we can assume $P(-M)=0$ and $P(M)=1$ for some large enough
$M$ (to be rigorous, we do the proof below for finite $M$ and a
RHS $\epsilon(M)$ instead of zero, then take the limit $M\rightarrow\infty$).
Since $\E_{x\sim p}x=0$, we have that
\[
\int_{-M}^{M}t\frac{dP(t)}{dt}\,dt=0
\]
Integration by parts gives that $\int_{-M}^{M}P(t)dt=M$. Note that
$P(t)$ is increasing logconcave, if $P(0)$ is too small, it would
make $\int_{-M}^{M}P(t)dt$ too small. To be precise, since $P$ is
logconcave, i.e., $-\log P(t)$ is convex, and so we have,
\[
-\log P(t)\geq-\log P(0)-\frac{\frac{d}{dt}P(0)}{P(0)}t.
\]
Or we simply write $P(t)\leq P(0)e^{\alpha t}$ for some $\alpha$.
Hence,
\[
M=\int_{-M}^{M}P(t)dt\leq\int_{-\infty}^{\frac{1}{\alpha}}P(0)e^{\alpha t}dt+\int_{1/\alpha}^{M}1dt=\frac{eP(0)}{\alpha}+M-\frac{1}{\alpha}.
\]
This shows that $P(0)\geq\frac{1}{e}$.

Next, Lemma \ref{lem:logcon_max-1} shows that $\max_{x}p(x)\le1$.
Therefore, the cumulative distribution $P$ is $1$-Lipschitz and
we have
\[
\P_{x\sim p}(x^{\top}\theta\geq t)\ge\P_{x\sim p}(x^{\top}\theta\geq0)-t\geq\frac{1}{e}-t.
\]
\end{proof}
\begin{lem}
\label{lem:logcon_max-1}Let $p$ be a one-dimensional isotropic logconcave
density$.$ Then $\max p(x)\le1$.
\end{lem}

For a proof of this (and for other properties of logconcave functions),
we refer the reader to \cite{lovasz2007geometry}.
\begin{xca}
Give a short proof that $\max_{x}p(x)=O(1)$ for any one-dimensional
isotropic logconcave density.
\end{xca}

Using the robust Gr\"unbaum theorem \ref{thm:grun-1}, we get the
following algorithm, which uses uniform random points from the current
set. Obtaining such a random sample algorithmically is an interesting
problem that we will study in the second part of this book.
\begin{lem}
Suppose $y^{(1)},\ldots,y^{(N)}$ are i.i.d. uniform random points
from a convex body $K$ and $y=\frac{1}{N}\sum_{i=1}^{N}y^{(i)}$.
Then for any halfspace $H$ not containing $y$, 
\[
\E(\vol(K\cap H))\le\left(1-\frac{1}{e}+\sqrt{\frac{n}{N}}\right)\vol(K).
\]
\end{lem}

\begin{proof}
Without loss of generality, we assume that $K$ is in isotropic position,
i.e., $\E_{K}(y^{(i)})=0$ and $\E_{K}(y^{(i)}(y^{(i)})^{\top})=I$.
Then we have $\E(y)=0$ and
\begin{align*}
\E\norm y^{2} & =\frac{1}{N^{2}}\E\sum_{i=1}^{N}\norm{y^{(i)}}^{2}=\frac{1}{N}\E\norm{y^{(1)}}^{2}=\frac{n}{N}.
\end{align*}
Therefore, by Jensen's inequality,
\[
\E\norm y\le\sqrt{\E\left(\norm y^{2}\right)}=\sqrt{\frac{n}{N}}.
\]
Thus, we can apply Theorem \ref{thm:grun-1} with $t=\sqrt{\frac{n}{N}}$
to bound $\E(\vol(K\cap H))/\vol(K)$. 
\end{proof}
Theorem \ref{thm:conv_opt} readily gives the following guarantee
for convex optimization, again using volume radius as the measure
of progress.
\begin{thm}
Let $f$ be a convex function on $\Rn$, $E^{(0)}$ be any initial
set and $\Omega\subset E^{(0)}$ be any convex set. Suppose that for
any $x\in E^{(0)}$, we can find a nonzero vector $g(x)$ such that
\[
f(y)\geq f(x)\text{ for any }y\text{ such that }g(x)^{\top}(y-x)\geq0.
\]

Then, for the center of gravity method with $N=10n$, we have
\[
\E\left(\min_{i=1,2,\cdots k}f(x^{(i)})\right)-\min_{y\in\Omega}f(y)\leq\left(\frac{\vol(E^{(0)})}{\vol(\Omega)}\right)^{\frac{1}{n}}\left(0.95\right)^{\frac{k}{n}}\left(\max_{z\in\Omega}f(z)-\min_{x\in\Omega}f(x)\right)\enspace.
\]
\end{thm}

Now we give a geometric proof of Theorem \ref{thm:Grunbaum}. Note
that one can modify the proof of Theorem \ref{thm:grun-1} to get
another proof.
\begin{proof}
Since affine transformations do not affect ratios of volumes, without
loss of generality, assume that the center of gravity of $K$ is the
origin and $H$ is the halfspace $\left\{ x:\:x_{1}\le0\right\} $.
For each point $t\in\R$, let $A(t)=K\cap\left\{ x:\ x_{1}=t\right\} $
be the $(n-1)$-dimensional slice of $K$ with $x_{1}=t$. Define
$r(t)$ as the radius of the $(n-1)$-dimensional ball with the same
$(n-1)$-dimensional volume as $A(t)$.

\begin{figure}[h]

\centering{}\includegraphics[width=0.85\columnwidth]{Figures/Thrm_3\lyxdot 20}\caption{Affine transformations for Centre of Gravity method}
\end{figure}

The goal of the proof is to show that the smallest possible halfspace
volume is achieved for a cone by a cut perpendicular to its axis.
In the first step, we will symmetrize $K$ as follows: replace each
cross-section $A(t)$ by a ball of the same volume, centered at $(t,0,\ldots,0)^{T}$.
We claim that the resulting rotationally symmetric body is convex.
To see this, note that all we have to show is that the radius function
$r(t)$ is concave. For any $s,t\in\R,$ and any $\lambda\in[0,1]$,
we have by convexity of $K$ that
\[
\lambda A(s)+(1-\lambda)A(t)\subseteq A(\lambda s+(1-\lambda)t)
\]
and by the Brunn-Minkowski theorem (Theorem \ref{thm:Brunn-Minkowski})
applied to $A(s),A(t)$, the function $\vol_{n-1}(A(s))^{1/(n-1)}$
is a concave function and so we have
\[
\vol_{n-1}(A(\lambda s+(1-\lambda)t))^{\frac{1}{n-1}}\ge\lambda\vol_{n-1}(A(s))^{\frac{1}{n-1}}+(1-\lambda)\vol_{n-1}(A(t))^{\frac{1}{n-1}}.
\]
From this and the definition of $r(t)$, it follows that
\[
r(\lambda s+(1-\lambda)t)\ge\lambda r(s)+(1-\lambda)r(t)
\]
as desired.

Next consider the subset $K_{1}=K\cap\left\{ x:\:x_{1}\le0\right\} $.
We replace this subset with a cone $C$ having the same base $A(0)$
and apex at some point along the $e_{1}$ axis so that the volume
of the cone is the same as $\vol(K_{1}).$ Using the concavity of
the radial function, this transformation can only decrease the center
of gravity along $e_{1}$. Therefore, proving a lower bound on the
transformed body $K_{1}$ will give a lower bound for $K$. So assume
we do this and the center of gravity is the origin. Next, extend the
cone to the right, so that it remains a rotational cone, and the volume
in the positive halfspace along $e_{1}$ is the same as $\vol(K\setminus K_{1})$.
Once again, the center of gravity can only move to the left, and so
the volume of $K_{1}$ can only decrease by this transformation. At
the end we have shown that the lower bound for any convex body follows
by proving for a rotational cone with axis along the normal to the
halfspace. The intersection of this cone with the halfspace is the
original cone scaled down by the ratio of the distance from the apex
to the center of gravity, to the height of the cone. So all that remains
to be done is to compute the relative distance of the center of gravity
from the apex, which is exactly $n/(n+1)$. Now we compute the volume
ratio:
\[
\frac{\vol(K_{1})}{\vol(K)}=\left(\frac{n}{n+1}\right)^{n}.
\]
\end{proof}
\begin{xca}
Show that for a cone of height $h$ in $\R^{n}$, i.e., the convex
hull of a convex body $K\subset\R^{n-1}$ with a single point $a\in\R^{n}$
at distance $h$ from the hyperplane $H$ containing $K$, the center
of gravity $z$ of the cone is at distance $h/(n+1)$ from $H$.
\end{xca}

To conclude this section, we note that the number of separation oracle
queries made by the center-of-gravity cutting plane method is asymptotically
the best possible.
\begin{thm}
Any algorithm that solves Problem \ref{prob:abstract_convex_problem}
using a separation oracle needs to make $\Omega(n\log(R/\epsilon))$
queries to the oracle.
\end{thm}

\begin{proof}
Suppose $K$ is a cube of side length $\epsilon$ contained in the
cube $[0,R]^{n}$. Imagine a tiling of the big cube by cubes of side
length $\epsilon$. Consider the oracle that always returns an axis
parallel halfspace that does not cut any little cube and contains
at least half of the volume of the ``remaining'' region, i.e., the
set given by the original cube intersected with all halfspaces given
by the oracle so far. This is always possible since for any halfspace
either the halfspace or its complement will contain at least half
the volume of any set. Thus each query at best halves the remaining
volume. To solve the problem, the algorithm needs to cut down to a
set of volume $\epsilon^{n}$ starting from a set of volume $R^{n}.$
Thus it needs at least $n\log_{2}(R/\epsilon)$ queries.
\end{proof}
\begin{figure}[h]
\centering{}\includegraphics[width=0.75\columnwidth]{Figures/Thrm_3\lyxdot 22}\caption{Complexity using a separation oracle}
\end{figure}


\section*{Discussion}

In later chapters we will see how to implement each iteration of the
center of gravity method in polynomial time. Computing the exact center
of gravity is \#P-hard even for a polytope \cite{rademacher2007approximating},
but we can find an arbitrarily close approximation in randomized polynomial
time via sampling. The method generalizes to certain noisy computation
models, e.g., when the oracle reports a function value that is within
a bounded additive error, i.e., a noisy function oracle.
